\documentclass[10pt,a4paper]{amsart}
\usepackage[margin=19mm]{geometry}
\usepackage{amsmath,amssymb,mathtools}
\usepackage[scr=rsfs,cal=euler]{mathalfa}
\usepackage{xfrac}
\usepackage[unicode,hidelinks,bookmarks=false]{hyperref}
\newcommand{\ie}{i.e.\ }
\newcommand{\cf}{cf.\ }
\newcommand{\resp}{respectively\ }
\newcommand{\Subsection}{\subsection}
\providecommand{\nobreakdash}{\nobreak\hbox{-}}
\setlength{\emergencystretch}{2em}
\multlinegap0pt
\usepackage{bm}
\usepackage{bbm}
\usepackage{dsfont}
\usepackage{xspace}
\usepackage{cancel}
\usepackage{mathtools}
\usepackage{amsthm}
\makeatletter
\@ifundefined{theorem}{\newtheorem{theorem}{Theorem}}{}
\makeatother
\title[A Distributional Approach to Generalized Stochastic Processe]{A Distributional Approach to Generalized Stochastic Processes on Locally Compact Abelian Groups}
\author{Hans G. Feichtinger, Wolfgang H\"ormann}
\date{Source: arXiv:2606.31316v1 (CC BY 4.0)}
\begin{document}
\maketitle

\noindent\emph{Source and licence.} A Distributional Approach to Generalized Stochastic Processes on Locally Compact Abelian Groups. Version arXiv:2606.31316v1 (CC BY 4.0), distributed under the Creative Commons Attribution 4.0 International licence, as recorded in the arXiv licence metadata for this exact version and shown on the abstract page https://arxiv.org/abs/2606.31316v1. Full rights notice: licenses/arxiv-2606.31316-CC-BY-4.0.txt.\par\smallskip

\noindent\textit{Editorial scope.} Counted unit: the Theorem whose source label is \texttt{ch18:the10} in the article (source0.tex lines 1264--1302, source label \texttt{ch18:the10}), delivered together with the article's own complete original proof of it (source0.tex lines 1304--1496, one \texttt{proof} environment of 193 lines). No mathematics is written, repaired or re-derived: statement and proof are the article's own text line for line. Required context retained uncounted: none. Every definition, display and label of those lines is retained unchanged. Omitted from this package and not redistributed: 1--1263, 1303--1303, 1497--2214. Nothing inside a retained excerpt is omitted or altered: every retained body is delivered exactly as the article's lines are written, so no omission receipt of any kind accompanies this package. Citations: the retained excerpts carry no citation command at all, so no bibliography entry of the article is transcribed and no reference is redistributed. Reference adaptation: none was needed and none was made; every reference inside the delivered unit resolves inside it, so no unresolved reference or citation is produced. Printed slips kept literal: no printed word, symbol, hypothesis or number of the counted statement or of its proof was altered. Layout and class: the article's own class is \texttt{svmult} with packages including graphicx, multicol, multirow, dcolumn, rotating, amssymb, amsfonts, amsmath, amscd, amsbsy, mathrsfs, bbm, subfigure, fontenc; the wrapper is the lane's pinned \texttt{amsart} editorial wrapper at 10pt on A4, which restates the theorem-like environments the retained text needs and adds the article's own macro definitions that the retained text needs (none), with extra wrapper packages (bm, bbm, dsfont, xspace, cancel, mathtools). The delivered numbering is the unit's own. Assets: no retained span contains a figure environment, a graphics-inclusion command, a PDF-page inclusion, a table or a TikZ picture; no image file is copied into this package, the package declares no asset dependency and no image file is redistributed with it. Licence: version arXiv:2606.31316v1 of this article is distributed under the Creative Commons Attribution 4.0 International licence, as recorded in the arXiv licence metadata for this exact version and shown on the abstract page https://arxiv.org/abs/2606.31316v1. A full rights notice is delivered at licenses/arxiv-2606.31316-CC-BY-4.0.txt. Characters: every retained excerpt is pure ASCII, so the delivered engine, XeLaTeX with its default Latin Modern text font, reports no missing character.
\begin{theorem}\label{ch18:the10}
Let $\rho $ be a GSP with covariance ${\sigma_{ \rho}}$:
\begin{enumerate}\leftskip5pt
\item[(a)] If $\sigma _{\rho} \in  S'_{0}(G\times G)
\text{is represented by some }
h \in  C^{b}(G\times G)$, then
 $\rho (f_{\alpha })$ is a Cauchy net in $\mathcal{H}$ whenever
$(f_{\alpha })_{\alpha \in I}$ is a vaguely convergent, $L^{1}$-bounded, and tight net in ${S_{0}(G)}$.

\item[(b)] In the above situation $\rho $ extends to
 a bounded linear operator
$\tilde{\rho}: M(G) \mapsto \mathcal{H}$ , which is $\sigma $-norm
continuous on tight subsets of $M(G)$, meaning that it converts
vaguely convergent, bounded, and  tight nets $(\mu_\alpha)$
into norm convergent nets $\tilde{\rho}(\mu_\alpha)$ in $\mathcal{H}$.
In particular,
$\lim_{y\rightarrow x} \tilde{\rho}(\delta _{y}) = \tilde{\rho}(\delta
_{x})$ in $\mathcal{H}$.  If $\left\{ \rho (f), f \in {S_{0}(G)} \right\} $ is dense in $\mathcal{H}$,
the extension $\tilde{\rho}$ is uniquely determined.

\item[(c)] The mapping $\rho_G \colon x \mapsto \tilde{\rho}(\delta_x)$ is a bounded, continuous stochastic process on $G$ and
$h(x,y) := (\tilde{\rho}(\delta_x) | \tilde{\rho}(\delta_y))$
is the covariance function of $\rho_G$.

\item[(d)] For any continuous and bounded stochastic process
 $\rho_{1}:G \mapsto  \mathcal{H}$,
the covariance function of $\rho _{1}  $ given by
$h(x,y) := (\rho_{1}(x)| \rho _{1}(y))$
is bounded and continuous on $G\times $G.  By vector-valued
integration, $\rho _{1}$ may be lifted to a bounded linear mapping
$\tilde{\rho}_{1}: M(G) \mapsto  \mathcal{H},$ which is $\sigma$-norm
continuous on bounded tight subsets. By
way of restriction to ${S_{0}(G)}$, $\tilde{\rho}_{1}$ may be considered
as a GSP  and $h$
represents the covariance distribution of this GSP.
%$\tilde{\rho}_{1}| _{\SOG}$.
% (by viewing $\Csp_b(G)$ as a subspace of $\SOPG$).
\end{enumerate}
\end{theorem}

\begin{proof}\
\begin{enumerate}\leftskip5pt
\item[(a)] Assume now that $(f_\alpha)$ is a net in ${S_{0}(G)}$, which  is a vaguely
convergent, bounded, and tight in $M(G)$.
To prove that the net
$\rho(f_{\alpha})_{\alpha \in I}$ is a Cauchy net
in $\mathcal{H}$, we use the fact that for any $k \in \mathcal{C}_c(G)$
one has
\begin{gather*}
\| \rho (f_{\alpha }){-}\rho (f_{\beta })\| ^{2}_{\mathcal{H}}
{=} (\rho (f_{\alpha }-f_{\beta })| \rho (f_{\alpha }-f_{\beta })) {=}
\langle\sigma _{\rho},(f_{\alpha }{-}f_{\beta })\otimes (\bar{f}_{\alpha
}{-}\bar{f}_{\beta })\rangle {=} \\
{=} \langle\sigma _{\rho} \cdot k \otimes k,(f_{\alpha }{-}f_{\beta
})\otimes
(\bar{f}_{\alpha }-\bar{f}_{\beta })\rangle + \langle\sigma _{\rho},k(f_{\alpha }-f_{\beta })\otimes (1-k)(\bar{f}_{\alpha
}-\bar{f}_{\beta })\rangle + \\
+ \langle\sigma _{\rho},(1-k)(f_{\alpha }-f_{\beta })\otimes
(\bar{f}_{\alpha }-\bar{f}_{\beta })\rangle.
\end{gather*}
As $f_{\alpha}$ is tight and $L^{1}$-bounded,
there exists $k \in  \mathcal{C}_c(G)$ with
$\| (1-k)$\break $(f_{\alpha } - f_{\beta })
\| _{1}< \epsilon  \ \forall \,  \alpha  \in I$. Thus
the second of the three terms can be handled in the following way:
\begin{gather*}
| \langle\sigma _{\rho},k(f_{\alpha }-f_{\beta })\otimes
(1-k)(\bar{f}_{\alpha }-\bar{f}_{\beta })\rangle| =
\\
{=} \left| \int^{}_{G} \left( \int^{}_{G}
h_{\rho}(x,y)k(x)(f_{\alpha }(x){-}f_{\beta }(x)) dx \right) (1{-}k(y))(\bar{f}_{\alpha }(y){-}\bar{f}_{\beta
}(y)) dy \right| \le
\\
\le  \| h_{\rho}\| _{\infty
} \int^{}_{G}    \left( \int^{}_{G} | k(x)(f_{\alpha }(x){-}f_{\beta }(x))|  dx \right) |
(1{-}k(y))(\bar{f}_{\alpha }(y)-\bar{f}_{\beta }(y))|  dy \le \\
\le  \| h_{\rho}\| _{\infty } \, C \int^{}_{G} |
(1-k(y))(\bar{f}_{\alpha }(y)-\bar{f}_{\beta }(y))|  dy
\end{gather*}
and the last integral can be made arbitrarily small by suitable
choice of $k$ above, \textit{independent} of $\alpha$ and $\beta$.
The third term can be treated in the same way as the second.
As the first term  converges to  0 if $f_{\alpha}$
is  vaguely convergent, it follows that
$\| \rho (f_{\alpha}) - \rho (f_{\beta})\| _{\mathcal{H}}$
tends to zero and this implies
that $(\rho (f_{\alpha }))_{\alpha \in I}$
is a Cauchy net in the norm topology.

\item[(b)] As any measure in $M(G)$ can be represented as the
$w^*$-limit of a bounded, tight net of functions in $S_0(G)$,
the density  of
$\left\{  \rho (f) , f \in {S_{0}(G)} \right\} $
in $\mathcal{H}$ implies the existence of a uniquely
determined element denoted by
$\tilde{\rho}(\mu ) \in  \mathcal{H}$,
such that $(\tilde{\rho}(\mu )| \rho (g)) := \lim_{\alpha }(\rho
(f_{\alpha })| \rho (g))  \ \forall \,  g \in {S_{0}(G)}$. The notation is
justified because it is independent of the choice of the net
$(f_{\alpha })$ with $\lim_{\alpha} f_{\alpha} = \mu$.
This can be seen using the proof of part a)
with two different $L^1$-bounded, tight nets with the
same vague limit.

Now we show the continuity of $\tilde{\rho}:$
%(\cite{HOERMD83}) \par
Let $(\mu _{\beta })_{\beta \in J}$  be a bounded and tight
net in $M(G),\,  w^\ast$-convergent  with  limit $\mu$.
Then we find that for any ${w}^*$-neighborhood
$  U := U(k_{1},k_{2},\ldots ,k_{n},\epsilon ) \in  \mathcal{U},$
 with $ k_{i} \in  {S_{0}(G)}  $, one can find for every $\beta$  some  $ f_{(U,\beta)} \in  {S_{0}(G)}$
such that the following holds:
\begin{gather*}
| \langle\mu _{\beta },k_{i}\rangle - \langle f_{(U,\beta)},k_{i}\rangle|
< \epsilon /2, \quad \forall \,  i=1\ldots n, \ \forall \, \beta\\
\text{and }\quad \quad | \langle\mu _{\beta },k_{i}\rangle - \langle\mu
,k_{i}\rangle|  < \epsilon /2,  \quad \forall \; \beta  > \beta _{0}.
\end{gather*}
Combining these two estimates one obtains
\[
| \langle \mu ,k_{i}\rangle - \langle
f_{(U,\beta)},k_{i}\rangle|  <  \epsilon,
 \quad \forall \,  i=1\ldots n \ \forall \,  \beta  > \beta _{0}.
\]
We have shown: $(f_{(U,\beta)})$ is a
$w^\ast$-convergent, $L^{1}$-bounded, and
tight net with the same limit
$\mu $ and $f_{(U,\beta)} \in {S_{0}(G)}$, for any
$U \in \mathcal{U}$.
Part a) implies that $\rho(f_{(U,\beta)})$ converges in
norm and
the limit is called $\tilde{\rho}(\mu )$ according to the above
definition.  Due to the construction of $f_{(U,\beta)}$, it is easy to
see that
$\tilde{\rho}(\mu _{\beta })$ converges to $\tilde{\rho}(\mu )$ as
well and this shows the ${w}^*$-to-norm continuity of $\tilde{\rho}$
on bounded, tight subsets. Furthermore, this implies the norm-norm
continuity and thus the boundedness of the mapping $\tilde{\rho}$
follows.

\item[(c)] The continuity of
$\rho _{G}$ follows from the ${w}^*$-to-norm continuity of
$\tilde{\rho}$ on tight subsets.
The boundedness of $\rho _{G}$ follows
from the fact that $\tilde{\rho}$ is bounded with respect to $\| .
\|_M$ together with
$\| \delta_x \|_M \,=\,1 \, \ \forall \, \, x \in G$.

 Let $(f_{\alpha })_{\alpha \in I}$  be tight, $L^{1}$-bounded
and  vaguely  convergent  with
limit $\delta _{0}$ (i.e., a generalized ``Dirac sequence'');
then the following completes the proof of part~c):
\begin{gather*}
h(x,y) = \int^{}_{G\times G} h(t) (\delta _{x}\otimes \delta _{y}) dt =
\lim_{\alpha } \langle h , T_{x}f_{\alpha }\otimes T_{y}f_{\alpha }
\rangle =\\
= \lim_{\alpha } \langle \sigma _{\rho} , T_{x}f_{\alpha }\otimes
T_{y}f_{\alpha } \rangle = (\tilde{\rho}(\delta _{x})| \tilde{\rho}(\bar{\delta }_{y})) = (\tilde{\rho}(\delta _{x})| \tilde{\rho}(\delta _{y})).
\end{gather*}

\item[(d)] The estimate
$ |h(x,y)| = |(\rho_{1}(x)| \rho _{1}(y))| \leq %hgfei
\| \rho_{1}(x)\|_{\mathcal{H}}\|\rho_{1}(y)\|_{\mathcal{H}} \le  c^{2}$
proves that $h$ is bounded. The continuity of $h$ results from the
continuity of $\rho _{1}$ and the inner product.
\end{enumerate}

With the help of vector-valued integration we can define
$ \tilde{\rho}_{1}(\mu)$ in the weak sense. In fact, for any
$ l \in \mathcal{H}$, the function $ (\rho _{1}(x)| l)$ is bounded and continuous;
hence the following integral makes sense:
\[
\bigl( \tilde{\rho}_{1}(\mu )| l\bigr) :=
\int^{}_{G} (\rho_{1}(x)| l) d\mu \quad \text{for} \,\, l \in \mathcal{H}\text{ and }\mu \in M(G).
\]
In view of the  Riesz  representation  theorem, $\tilde{\rho}_{1}(\mu )$  is  a  well-defined element of $\mathcal{H}$,
as $|(\tilde{\rho}_1(\mu)|l)| \: \leq \: c  \| \mu \|_M
\| l \|_{\mathcal{H}}$. The last inequality implies
$\| \tilde{\rho}_1(\mu) \|_\mathcal{H} \: \leq \: c \|
\mu \|_M$ and thus the boundedness of $\widetilde{\rho_1}$ with
respect to $\| . \|_M$.



To prove the ${w}^*$-to-norm continuity, we take a bounded, tight,
$w^\ast$-convergent net $\mu_{\alpha}$ in $M(G)$ with limit $\mu$. As
the mapping $x \mapsto ({\rho_1}(x)|l)$ is continuous and bounded for
any $l \in \mathcal{H}$ and $\mu_{\alpha}$ is tight we get it is enough
to know that the integrand is uniformly bounded and uniformly convergent
over compact subset of $G$, in order to conclude that
%
\begin{align*}
\lim_{\alpha}(\tilde{\rho}_1(\mu_{\alpha})|l) & = \lim_{\alpha} \int_G({\rho}_1(x)|l) d\mu_{\alpha}\\
& = \int_G({\rho}_1(x)|l) d\mu
 = (\tilde{\rho}_1(\mu)|l),
\end{align*}
%
which shows the ${w}^*-{w}^*$ continuity of $\tilde{\rho}_1$.
The convergence of
$\|\tilde{\rho}_1(\mu_{\alpha}) \|$
is shown by the following equality:
%
\begin{align*}
\lim_{\alpha} (\tilde{\rho}_1(\mu_{\alpha}) | \tilde{\rho}_1(\mu_{\alpha})) & = \lim_{\alpha }\int_G \int_G
(\rho_1(x)|\rho_1(y)) d\mu_{\alpha} d\mu_{\alpha}\\
& = \int_G \int_G (\rho_1(x)|\rho_1(y)) d\mu d\mu \quad
  = \, \,  \|\tilde{\rho}_1(\mu)\|^2,
\end{align*}
 %
which is true as $(\rho_1(x)|\rho_1(y))$ is continuous and
bounded and $\mu_{\alpha}$ is $w^\ast$-convergent,
bounded, and tight. The last result together with the
${w}^*$-to-norm continuity of $\tilde{\rho}_1$
implies the ${w}^*$-to-norm continuity, and our claim is proved.
Furthermore %hgfei 
we get the required GSP by restriction to $ {S_{0}(G)}$:
$\rho  := \tilde{\rho}_{1}| _{S_{0}}$.


It remains to be shown that
$h(x,y) = (\rho _{1}(x)| \rho_{1}(y))$
represents  the covariance distribution
$\sigma _{\rho}$ of $\rho $. This follows from the identity
\begin{gather*}
\langle \sigma _{\rho},f\otimes g \rangle = \bigl( \rho (f)| \rho
(\bar{g})\bigr)  = \int^{}_{G} (\rho _{1}(x)| \rho (\bar{g})) f(x) dx \\ \quad\quad\quad\quad\quad\quad\quad=
\int^{}_{G} \int^{}_{G} g(y)(\rho _{1}(x)| \rho _{1}(y)) dy f(x) dx\\
=
\int^{}_{G} \int^{}_{G}h(x,y) f(x) g(y) dx dy
= \langle h,f\otimes g\rangle \text{ for }f,g \in {S_{0}(G)}.
\quad \quad \qed 
\end{gather*}
\end{proof}

\end{document}
